% реферат по теме 'АИС и АСУ'
\documentclass{SibFU-docs}

% доп. пакеты
\usepackage{graphicx}
\usepackage{float}
\usepackage{hyperref}
\usepackage{caption}
\usepackage{subcaption}
\usepackage{amsmath}
% предотвартим ошибку с \Bbbk перед загрузкой пакета amssymb
\let\Bbbk\relax
\usepackage{amssymb}
\usepackage{booktabs}
\usepackage{tabularx}
\usepackage{longtable}
\usepackage{enumitem}
\usepackage{multirow}

% гиперссылки
\hypersetup{
    colorlinks=true,
    linkcolor=black,
    citecolor=blue,
    urlcolor=blue
}

% библиография
\addbibresource{report.bib}

% титульный лист
\begin{document}
\setlist[itemize]{left=1.3cm, itemsep=1pt, topsep=2pt} % для выравнивания списков
\setlist[enumerate]{left=1.3cm, itemsep=1pt, topsep=2pt}
\renewcommand{\contentsname}{}
\mycovertitle
{Гуманитарный институт}
{Кафедра прикладной информатики в искусстве и интерактивном медиа}
{Автоматизированные информационные системы и автоматизированные системы управления}
{ст. преподаватель И.~Р.~Нигматуллин}
{Логинова Ю.В., Захаров И.М.}
 
\begin{center}
\bfseries РЕФЕРАТ
\end{center}
\noindent

Реферат по теме: «Автоматизированные информационные системы и автоматизированные системы управления» содержит 30 страниц и 18 использованных источников.

Ключевые слова: капсом, основное, сплошняком, через запятую
АВТОМАТИЗИРОВАННЫЕ ИНФОРМАЦИОННЫЕ СИСТЕМЫ, АВТОМАТИЗИРОВАННЫЕ СИСТЕМЫ УПРАВЛЕНИЯ, НАЗНАЧЕНИЕ, ФУНКЦИИ, АРХИТЕКТУРА, КЛЮЧЕВЫЕ КОМПОНЕНТЫ, ЖИЗНЕННЫЙ ЦИКЛ, ЭТАПЫ РАЗРАБОТКИ, СОВРЕМЕННЫЕ ТЕХНОЛОГИИ, ПРОГРАММНОЕ ОБЕСПЕЧЕНИЕ, БЕЗОПАСНОСТЬ, УСТОЙЧИВОСТЬ, ИНФОРМАЦИОННАЯ ЗАЩИТА, КИБЕРБЕЗОПАСНОСТЬ, ПРИМЕНЕНИЕ, ПРОМЫШЛЕННОСТЬ, БИЗНЕС, ГОСУДАРСТВЕННОЕ УПРАВЛЕНИЕ, ЭФФЕКТИВНОСТЬ, УПРАВЛЕНЧЕСКИЕ РЕШЕНИЯ.

В работе исследуются сущность, архитектура и практическое применение автоматизированных информационных систем (АИС) и автоматизированных систем управления (АСУ). Рассмотрены их ключевые характеристики, структурные компоненты, этапы жизненного цикла, а также современные технологические решения для их реализации. Особое внимание уделено вопросам обеспечения безопасности, устойчивости функционирования и конкретным примерам внедрения в различных отраслях.

Основные выводы:

\begin{enumerate}
\item Автоматизированные информационные системы и системы управления являются технологическим и организационным ядром цифровой трансформации, кардинально повышая эффективность обработки данных и качество управленческих решений за счёт формализации и алгоритмизации бизнес-процессов.

\item Архитектура АСУ, основанная на чётком взаимодействии компонентов (информационного, программного, технического, организационного), определяет её надёжность, гибкость и способность к интеграции с внешними сервисами и системами.

\item Строгое следование этапам жизненного цикла (проектирование, разработка, внедрение, эксплуатация, сопровождение, модернизация) является критическим фактором успеха создания и долгосрочной эффективности автоматизированных систем, минимизируя риски и затраты.

\item Развитие современных технологий (облачные сервисы, IoT, Big Data, AI) трансформирует АИС и АСУ в интеллектуальные платформы, способные не только автоматизировать рутинные операции, но и обеспечивать прогнозную аналитику и адаптивное управление в реальном времени.
\end{enumerate}

Рекомендации:

\begin{itemize}[label=-]
\item При проектировании цифровой инфраструктуры организации рассматривать ИТ и ИС как единый комплекс, где технологии служат реализации системных целей.
\item Внедрять ИС, соответствующие уровню зрелости бизнес-процессов организации — от автоматизации транзакций (OLTP) до аналитической обработки (OLAP) и интеллектуального анализа данных.
\item Формировать корпоративные хранилища данных (Data Warehouse) как фундамент для эффективного управления информационными ресурсами и аналитики.
\item Использовать возможности современных ИС (ERP, CRM, BI-платформы) для создания единого информационного пространства и снижения рисков при принятии управленческих решений.
\end{itemize}

\newpage

\thispagestyle{empty}
\begin{center}
\bfseries СОДЕРЖАНИЕ
\end{center}
\vspace{-6pt}
\tableofcontents
\clearpage

% фантомные боли для секций, чтобы не было нумераций кроме тела реферата
\section*{}
\vspace*{-2.5em}
\begin{center}\textbf{ВВЕДЕНИЕ}\end{center}
\phantomsection
\addcontentsline{toc}{section}{Введение}

Автоматизированные информационные системы и системы управления превратились из вспомогательного инструментария в критическую инфраструктуру, определяющую эффективность и конкурентоспособность практически любой организации — от промышленного предприятия и финансовой корпорации до государственного учреждения. Их роль эволюционировала от простой механизации расчётов к созданию комплексных цифровых экосистем, управляющих потоками данных, ресурсами и процессами в режиме реального времени.

Современное состояние проблемы. Стремительное развитие вычислительных мощностей, сетевых технологий и алгоритмов искусственного интеллекта постоянно расширяет функциональные границы АИС и АСУ. Однако это приводит к возрастающей сложности их архитектуры, интеграционных задач и, как следствие, к новым вызовам в области кибербезопасности, надёжности и управления жизненным циклом. Разрыв между технологическими возможностями и умением организаций грамотно проектировать, внедрять и эффективно эксплуатировать такие системы остаётся значительным.

Актуальность темы обусловлена всеобщей цифровизацией, при которой корректное построение и эксплуатация АИС/АСУ напрямую влияют на операционную деятельность, стратегическое развитие и устойчивость организаций перед лицом внутренних и внешних угроз.

Новизна исследования заключается в комплексном рассмотрении АИС и АСУ как взаимосвязанных и эволюционирующих комплексов, с акцентом на системный подход к их архитектуре, жизненному циклу и практическим аспектам обеспечения безопасности и устойчивости в различных предметных областях.

Цель работы: исследовать теоретические основы, принципы построения, современное состояние и практику применения автоматизированных информационных систем и систем управления.

Задачи исследования:
раскрыть назначение, классификацию и основные характеристики автоматизированных информационных систем;
проанализировать архитектурные принципы и ключевые компоненты автоматизированных систем управления;
изучить этапы жизненного цикла и методологии разработки автоматизированных систем;
охарактеризовать современные технологии и типовые программные решения для автоматизации управления;
исследовать вопросы обеспечения безопасности, надёжности и отказоустойчивости автоматизированных систем;
привести и проанализировать примеры успешного применения АИС и АСУ в различных сферах деятельности.

Методы исследования:
системный анализ для изучения АИС и АСУ как целостных организационно-технических комплексов;
структурно-функциональный подход для исследования архитектуры и взаимодействия компонентов;
сравнительный анализ современных технологических платформ и программных решений;
анализ литературных источников, отраслевых стандартов и практических кейсов внедрения;
обобщение лучших практик в области управления жизненным циклом и обеспечения информационной безопасности систем.

\newpage
% вручную увеличить номер секции и записать в оглавление с номером,
% чтобы в теле не печатался автоматический номер, но в tableofcontents он остался
\refstepcounter{section}
\addcontentsline{toc}{section}{\protect\numberline{\thesection} Назначение и основные характеристики автоматизированных информационных систем}
\vspace*{-2.5em}
\begin{center}\textbf{\thesection\ НАЗНАЧЕНИЕ И ОСНОВНЫЕ ХАРАКТЕРИСТИКИ АВТОМАТИЗИРОВАННЫХ ИНФОРМАЦИОННЫХ СИСТЕМ}\end{center}
\vspace*{-1.5em}
\ManualSubsection{Предназначение автоматизированных информационных систем}

Автоматизированная информационная система (АИС) представляет собой комплекс программных, технических и организационных средств, предназначенный для автоматизации процессов сбора, хранения, обработки и предоставления информации в рамках конкретной предметной области. Ключевая цель внедрения АИС заключается в повышении эффективности работы организации за счет минимизации рутинных операций и снижения влияния человеческого фактора. Как отмечается в источниках, такие системы призваны упорядочить информационный хаос, взяв на себя техническую часть работы с данными, что позволяет сотрудникам сосредоточиться на решении содержательных задач \cite{ais_basics}.

Основными целями внедрения автоматизированных систем являются ускорение бизнес-процессов, повышение точности и достоверности информации, систематизация данных, автоматизация формирования отчетности и усиление контроля над операционной деятельностью. Внедрение АИС приводит к трансформации рабочих процессов, где задачи, требовавшие ранее значительных временных затрат, выполняются за счет нескольких кликов, а данные становятся доступными в режиме реального времени для принятия обоснованных управленческих решений \cite{ais_efficiency}.

\ManualSubsection{Ключевые характеристики автоматизированных информационных систем}

Эффективность автоматизированной информационной системы определяется набором ее базовых характеристик, которые обеспечивают практическую ценность для пользователей. Согласно общепринятым подходам, к таким характеристикам относятся полнота, точность, своевременность, доступность и защищенность \cite{ais_components}. Полнота подразумевает охват системой всех необходимых для деятельности организации процессов и функций, что исключает необходимость использования несвязанных между собой программных решений. Точность гарантирует актуальность и достоверность хранимой и обрабатываемой информации, что является фундаментом для доверия к системе и принятия на ее основе ответственных решений.

Своевременность обеспечивает предоставление информации без задержек, отражая изменения в данных сразу после их возникновения, что критически важно в динамичной бизнес-среде. Доступность означает простоту и скорость поиска необходимых сведений для авторизованных пользователей, устраняя временные потери на рутинный поиск. Защищенность является комплексным свойством, гарантирующим сохранность, конфиденциальность и целостность данных от несанкционированного доступа и утери. Совокупность этих характеристик преобразует АИС из простого инструмента учета в стратегический актив организации.

\ManualSubsection{Структурные компоненты автоматизированных систем}

Архитектура автоматизированной информационной системы представляет собой совокупность взаимосвязанных компонентов, каждый из которых выполняет строго определенную роль. Основу системы составляет техническое обеспечение, включающее компьютеры, серверы, сетевое оборудование и периферийные устройства, которые формируют физическую платформу для функционирования АИС. Программное обеспечение выступает в роли «мозга» системы, включая операционные системы, системы управления базами данных (СУБД) и специализированные прикладные программы, реализующие бизнес-логику \cite{ais_components}.

Информационное обеспечение представляет собой содержание системы — структурированные массивы данных, документы, классификаторы и базы знаний, с которыми непосредственно работает пользователь. Организационное обеспечение включает регламенты, инструкции, методологии и стандарты, которые определяют порядок взаимодействия пользователей с системой и между собой. Кадровое обеспечение является критически важным компонентом, так как охватывает администраторов, поддерживающих работоспособность системы, и конечных пользователей, которые применяют ее для решения профессиональных задач. Отсутствие или слабая проработка любого из этих компонентов сводит на нет эффективность всей системы.

\ManualSubsection{Области практического применения АИС}

Автоматизированные информационные системы нашли широкое применение практически во всех сферах экономической и социальной деятельности, адаптируясь к специфике различных предметных областей. В коммерческом секторе наиболее распространены системы управления взаимоотношениями с клиентами (CRM) и системы планирования ресурсов предприятия (ERP), которые автоматизируют учет, продажи, логистику и финансовое планирование. В производственной сфере используются автоматизированные системы управления технологическими процессами (АСУ ТП), осуществляющие контроль за оборудованием, параметрами производства и качеством выпускаемой продукции.

В области здравоохранения медицинские информационные системы (МИС) обеспечивают ведение электронных историй болезни, запись пациентов на прием и управление медицинскими ресурсами. В образовательных учреждениях применяются системы управления обучением (LMS), автоматизирующие составление расписаний, ведение электронных журналов и предоставление учебных материалов. Государственный сектор активно внедряет различные государственные информационные системы (ГИС), предназначенные для оказания услуг гражданам в электронном виде, межведомственного документооборота и управления. Это разнообразие применений демонстрирует универсальность концепции АИС как инструмента повышения эффективности управления данными и процессами.

\ManualSubsection{Преимущества внедрения и критерии выбора АИС}

Внедрение автоматизированной информационной системы приносит организации ряд значимых преимуществ, напрямую влияющих на ее конкурентоспособность и операционную эффективность. Ключевым из них является существенная экономия времени за счет автоматизации рутинных операций, что позволяет персоналу перераспределить усилия на решение аналитических и стратегических задач. Снижение количества ошибок, обусловленное минимизацией ручного ввода данных и автоматизацией расчетов, повышает надежность и качество результатов работы. Повышение уровня контроля обеспечивается за счет прозрачности всех фиксируемых в системе операций и наличия инструментов аудита \cite{ais_efficiency}.

При выборе конкретной АИС для организации необходимо учитывать несколько критически важных критериев. Первичным является соответствие системы стоящим перед организацией задачам и наличие в ней необходимого функционального покрытия. Простота освоения и удобство интерфейса напрямую определяют скорость адаптации сотрудников и, в конечном итоге, успешность внедрения. Масштабируемость архитектуры системы должна позволять ей расти вместе с бизнесом, поддерживая увеличение числа пользователей и объема данных. Не менее важны наличие и качество технической поддержки со стороны вендора, а также анализ совокупной стоимости владения, включающей не только первоначальные лицензионные расходы, но и затраты на обновление, обслуживание и обучение.

\newpage
\refstepcounter{section}
\addcontentsline{toc}{section}{\protect\numberline{\thesection}Архитектура и ключевые компоненты автоматизированных систем управления}
\begin{center}\textbf{\thesection\ АРХИТЕКТУРА И КЛЮЧЕВЫЕ КОМПОНЕНТЫ АВТОМАТИЗИРОВАННЫХ СИСТЕМ УПРАВЛЕНИЯ}\end{center}
\vspace*{-1.5em}
\ManualSubsection{Понятие и подходы к архитектуре систем управления}

Архитектура автоматизированной системы управления (АСУ) представляет собой фундаментальную организацию данной системы, воплощённую в её компонентах, их взаимосвязях, а также в принципах её проектирования и эволюции \cite{gosstandart}. В сущности, это схема, определяющая, как все элементы системы — от технических устройств до программных модулей — соединены и взаимодействуют для достижения общей цели управления. Архитектура служит связующим звеном между требованиями управленческих процессов и их технической реализацией, обеспечивая целостность и эффективность системы. В рамках современного подхода бизнес-архитектура фокусируется на отображении ключевых элементов организации и их взаимоотношений, что является основой для проектирования ИТ-решений \cite{iiba_components}.

Существует несколько типовых подходов к организации архитектуры, каждый из которых имеет определённую область применения. Централизованная архитектура, где управление сосредоточено в одном основном вычислительном узле, подходит для компактных систем с локализованными компонентами. Децентрализованная архитектура предполагает распределение функций управления между несколькими узлами, что характерно для крупных, территориально распределённых организаций. Клиент-серверная модель, в которой сервер обрабатывает данные, а клиенты предоставляют интерфейс, широко используется в современных офисных и бизнес-системах. Облачная архитектура, размещающая компоненты в интернет-пространстве, является основой для современных цифровых проектов и мобильных сервисов \cite{rasp_control_system}.

\ManualSubsection{Функциональные компоненты системы управления}

Любая автоматизированная система управления включает в себя ряд обязательных функциональных компонентов, которые в совокупности образуют замкнутый контур управления. Согласно описаниям и примерам, ядро системы составляют блок сбора информации, блок обработки данных, блок принятия решений, интерфейс пользователя и хранилище данных \cite{rasp_control_system}. Блок сбора информации отвечает за получение первичных данных от оборудования, датчиков или пользовательских интерфейсов, используя такие устройства, как сенсоры, приборы учета и сканеры.

Блок обработки данных выполняет анализ и преобразование собранной информации, реализуя алгоритмы фильтрации, агрегации и первичной аналитики, что обычно происходит на серверных станциях или специализированных компьютерах. Блок принятия решений формирует управляющие команды на основе обработанных данных и заложенных алгоритмов, используя контроллеры и специализированное программное обеспечение. Интерфейс пользователя обеспечивает визуализацию состояния системы, отображение данных и приём команд от оператора через экраны мониторов, панели управления или веб-сайты. Хранилище данных выполняет функцию долговременного сохранения как исходной информации, так и результатов её обработки, используя базы данных и системы хранения файлов.

\ManualSubsection{Принципы проектирования архитектуры АСУ}

Проектирование архитектуры автоматизированной системы управления основывается на ряде ключевых принципов, направленных на обеспечение её долгосрочной эффективности и жизнеспособности. Масштабируемость (возможность роста) является фундаментальным принципом, позволяющим наращивать функциональность и мощность системы по мере роста потребностей организации без её кардинальной перестройки. Надёжность и отказоустойчивость обеспечивают стабильность работы и непрерывность критических бизнес-процессов, что достигается за счёт резервирования компонентов и продуманного управления сбоями.

Безопасность (защита информации) представляет собой сквозной принцип, требующий реализации механизмов для охраны данных от несанкционированного доступа на всех уровнях архитектуры. Адаптивность подразумевает способность системы к относительно безболезненной модификации в ответ на изменения внешних условий или внутренних требований. Совместимость (способность к совместной работе) обеспечивает возможность интеграции системы с существующим оборудованием и программным обеспечением, защищая предыдущие технологические инвестиции организации. Соблюдение этих принципов в ходе проектирования напрямую влияет на стоимость владения, гибкость и конечную успешность внедряемого решения.

\ManualSubsection{Особенности архитектурных решений в различных отраслях}

Специфика различных отраслей накладывает отпечаток на выбор и реализацию архитектурных решений для автоматизированных систем управления. В промышленном производстве, как правило, применяется многоуровневая иерархическая архитектура. На нижнем, операционном уровне находятся датчики и программируемые логические контроллеры (ПЛК), непосредственно взаимодействующие с производственным оборудованием. Цеховой уровень объединяет данные с нескольких линий или участков, обеспечивая мониторинг и локальное управление. Корпоративный уровень предоставляет руководству агрегированную информацию и инструменты для стратегического планирования, интегрируя производственные данные с системами ERP.

В коммерческих организациях и сфере услуг доминирует клиент-серверная архитектура. Центральный сервер или их кластер выступает в роли единого источника истины, храня всю актуальную информацию. Сотрудники получают доступ к данным и функционалу через рабочие станции, ноутбуки или мобильные устройства, выступающие в роли клиентов. В системах автоматизации зданий (Building Management Systems, BMS) часто используется децентрализованный подход с элементами интеграции. Отдельные подсистемы управления освещением, отоплением, вентиляцией и безопасностью функционируют автономно, но координируются через единый диспетчерский центр, что повышает общую энергоэффективность и удобство управления \cite{rasp_control_system}.

\newpage
\refstepcounter{section}
\addcontentsline{toc}{section}{\protect\numberline{\thesection}Жизненный цикл и этапы разработки автоматизированных систем}
\vspace*{-1.5em}
\begin{center}\textbf{\thesection\ ЖИЗНЕННЫЙ ЦИКЛ И ЭТАПЫ РАЗРАБОТКИ АВТОМАТИЗИРОВАННЫХ СИСТЕМ}\end{center}
\vspace*{-1.5em}
\ManualSubsection{Концепция жизненного цикла автоматизированных систем}

Жизненный цикл автоматизированной системы (АС) представляет собой полную последовательность взаимосвязанных этапов её существования — от момента зарождения первоначальной концепции до окончательного вывода из эксплуатации и утилизации \cite{gost-34}. Это системный взгляд на АС как на развивающийся объект, имеющий определённые фазы развития, каждая из которых характеризуется специфическими целями, задачами и результатами. Понимание жизненного цикла имеет фундаментальное значение для эффективного управления проектами по созданию и сопровождению АС, так как позволяет структурировать работы, рационально распределять ресурсы и минимизировать риски на всех стадиях. Международный стандарт ISO/IEC 12207 предоставляет общую структуру процессов жизненного цикла программного обеспечения, которая может быть адаптирована для проектов автоматизированных систем \cite{iso-iec-12207}.

\ManualSubsection{Стадии предпроектного анализа и проектирования}

Начальные этапы жизненного цикла закладывают основу для успеха всего проекта. Стадия предпроектного анализа (идеи и планирования) посвящена формированию концепции системы. В её рамках проводится анализ существующих бизнес-процессов организации, выявляются потребности и требования будущих пользователей, оцениваются технико-экономические аспекты реализации. Ключевым результатом данного этапа является разработка технического задания (ТЗ), которое служит формальным соглашением между заказчиком и разработчиком, фиксирующим цели, задачи, основные функции и ограничения будущей системы \cite{swebok}.

Следующая стадия — техническое проектирование — трансформирует требования ТЗ в детализированные технические решения. На этом этапе определяется архитектура системы, проектируются интерфейсы, модели данных и алгоритмы, выбираются технологии и инструменты разработки. Создаются прототипы ключевых компонентов для валидации архитектурных решений с заказчиком. Итогом фазы проектирования является комплект проектной документации (технический проект), детально описывающий внутреннее устройство системы и служащий руководством для её непосредственной реализации.

\ManualSubsection{Стадии разработки, внедрения и эксплуатации}

Стадия разработки (реализации) — это процесс непосредственного создания компонентов системы в соответствии с утверждённым техническим проектом. Он включает в себя программирование, создание баз данных, разработку пользовательских интерфейсов, интеграцию отдельных модулей в единый комплекс. Параллельно с разработкой ведётся тестирование (unit-тесты, интеграционное тестирование), а также создаётся эксплуатационная документация \cite{iso-iec-12207}.

Стадия внедрения (ввода в действие) представляет собой комплекс мероприятий по установке созданной системы в производственную среду организации. Это включает установку и настройку программного и аппаратного обеспечения, миграцию исторических данных, обучение персонала и проведение опытной эксплуатации. Внедрение часто осуществляется поэтапно, чтобы минимизировать риски для текущей деятельности организации. После успешного завершения опытной эксплуатации система переходит в стадию промышленной эксплуатации. На этой стадии АС используется для решения реальных бизнес-задач. Основными процессами являются оперативное управление системой, поддержка пользователей, мониторинг производительности и надёжности, а также сбор обратной связи для последующего развития.

\ManualSubsection{Стадии сопровождения и завершения жизненного цикла}

Стадия сопровождения (поддержки и развития) является наиболее продолжительной в жизненном цикле. Её задача — обеспечить непрерывную и эффективную работу системы в условиях изменяющихся требований бизнеса и технологической среды. Сопровождение включает в себя исправление обнаруженных ошибок, адаптацию системы к изменениям в законодательстве или бизнес-процессах, а также её совершенствование путём добавления нового функционала. Модернизация системы на данном этапе позволяет продлить срок её полезного использования и повысить возврат инвестиций.

Завершающей фазой жизненного цикла является вывод системы из эксплуатации. Данный этап инициируется, когда дальнейшая поддержка и модернизация становятся экономически нецелесообразными или технически невозможными. Процесс вывода включает планирование остановки, архивирование всех исторических данных в соответствии с законодательными и корпоративными требованиями, подготовку отчётов о результатах эксплуатации, а также безопасное изъятие и утилизацию устаревшего оборудования. Грамотное завершение жизненного цикла важно для обеспечения преемственности и перехода на новую, более современную систему.

\newpage
\refstepcounter{section}
\addcontentsline{toc}{section}{\protect\numberline{\thesection}Современные технологии и программные решения для автоматизации управления}
\begin{center}\textbf{\thesection\ СОВРЕМЕННЫЕ ТЕХНОЛОГИИ И ПРОГРАММНЫЕ РЕШЕНИЯ ДЛЯ АВТОМАТИЗАЦИИ УПРАВЛЕНИЯ}\end{center}
\vspace*{-1.5em}
\ManualSubsection{Технологические драйверы современной автоматизации}

Современный этап развития автоматизированных систем управления характеризуется конвергенцией ряда технологических трендов, трансформирующих традиционные подходы. Ключевыми драйверами выступают облачные вычисления, аналитика больших данных (Big Data), искусственный интеллект (ИИ) и интернет вещей (IoT), которые в совокупности способствуют созданию более интеллектуальных, адаптивных и гибких систем \cite{gartner-2024}. Эти технологии позволяют системам перейти от простого выполнения регламентированных операций к аналитической поддержке принятия решений, прогнозированию и оптимизации процессов. Массовое внедрение облачных сервисов, предлагающих модели «программное обеспечение как услуга» (SaaS), сделало мощные инструменты управления доступными для организаций любого масштаба, устраняя высокие первоначальные затраты на инфраструктуру и обеспечивая гибкое масштабирование ресурсов \cite{deloitte-erp-cloud-2023}.

\ManualSubsection{Эволюция ключевых классов программных решений}

Ключевые классы программных решений для управления, такие как системы планирования ресурсов предприятия (ERP) и системы управления взаимоотношениями с клиентами (CRM), претерпели значительную эволюцию. Современные ERP-системы нового поколения, базирующиеся на облачных архитектурах, стали более гибкими, модульными и ориентированными на пользовательский опыт. Они предлагают глубокую интеграцию между финансовыми, операционными и кадровыми процессами, обеспечивая единый источник достоверных данных для всей организации. CRM-системы трансформировались из простых баз данных клиентов в комплексные платформы, которые автоматизируют полный цикл продаж и обслуживания, обогащаясь аналитическими модулями на основе ИИ для сегментации клиентов и прогнозирования их поведения \cite{mckinsey-automation-2023}.

Параллельно возникли и набрали популярность новые парадигмы разработки и автоматизации. Low-code/no-code платформы позволяют бизнес-пользователям создавать и модифицировать приложения с минимальным привлечением профессиональных разработчиков, используя визуальные конструкторы и готовые компоненты. Технологии роботизации процессов (Robotic Process Automation, RPA) предоставляют инструменты для создания программных роботов, автоматизирующих рутинные, повторяющиеся задачи по работе с данными в различных, зачастую неинтегрированных между собой, информационных системах.

\ManualSubsection{Интеграция аналитики и искусственного интеллекта}

Интеграция аналитики больших данных и методов искусственного интеллекта представляет собой качественный скачок в возможностях систем управления. Технологии машинного обучения позволяют системам выявлять сложные, неочевидные паттерны и зависимости в исторических данных, формируя основу для предиктивной аналитики. Это находит применение в прогнозировании спроса, оптимизации цепочек поставок, превентивном обслуживании оборудования и оценке рисков. Обработка естественного языка (NLP) лежит в основе интеллектуальных чат-ботов и виртуальных ассистентов, которые берут на себя первичное взаимодействие с клиентами и сотрудниками, обрабатывая запросы и извлекая структурированную информацию из неструктурированных текстовых данных \cite{gartner-2024}.

Распространение интернета вещей (IoT) обеспечивает новый уровень сбора операционных данных с физических активов — от промышленного оборудования до офисной инфраструктуры. Потоки данных с датчиков в реальном времени позволяют реализовать концепцию «цифрового двойника» (Digital Twin) — виртуальной модели физического объекта или процесса, которая используется для мониторинга, анализа и симуляции. Это создаёт основу для создания полностью автоматизированных, самооптимизирующихся систем управления производством, логистикой и инфраструктурой.

\ManualSubsection{Тенденции развития и практические аспекты внедрения}

Основными тенденциями развития на ближайшую перспективу являются углубление интеграции между различными технологическими стеками, рост влияния искусственного интеллекта принятия решений (Decision Intelligence) и повышение значимости пользовательского опыта (UX) при проектировании систем. Ускоренное развёртывание сетей пятого поколения (5G) будет способствовать дальнейшему распространению IoT-решений и мобильных сервисов, требующих высокой скорости передачи данных и низкой задержки \cite{gartner-2024}.

При выборе и внедрении современных технологий автоматизации организациям рекомендуется придерживаться стратегического подхода, ориентированного на решение конкретных бизнес-задач, а не на следование технологическим модам. Критически важным является оценка общей стоимости владения, включающей расходы на интеграцию с существующим ландшафтом информационных систем, обучение персонала, техническую поддержку и регулярные обновления. Поэтапный, итеративный подход к внедрению, начиная с автоматизации наиболее проблемных или высокозатратных процессов, позволяет быстрее получить измеримую отдачу от инвестиций и накопить необходимый опыт для реализации более масштабных преобразований.

\newpage
\refstepcounter{section}
\addcontentsline{toc}{section}{\protect\numberline{\thesection}Вопросы безопасности и устойчивости функционирования автоматизированных систем}
\begin{center}\textbf{\thesection\ ВОПРОСЫ БЕЗОПАСНОСТИ И УСТОЙЧИВОСТИ ФУНКЦИОНИРОВАНИЯ АВТОМАТИЗИРОВАННЫХ СИСТЕМ}\end{center}
\vspace*{-1.5em}
\ManualSubsection{Концептуальные основы информационной безопасности}

Информационная безопасность автоматизированных систем представляет собой комплексную дисциплину, направленную на обеспечение конфиденциальности, целостности и доступности данных и сервисов, обрабатываемых этими системами. В современных условиях, характеризующихся повсеместной цифровизацией и интернет-подключённостью, обеспечение безопасности перестаёт быть факультативной опцией и становится критическим требованием для устойчивого функционирования любой организации. Угрозы безопасности носят многофакторный характер и включают как целенаправленные внешние атаки (киберпреступность, целевые взломы), так и внутренние риски, связанные с действиями персонала, техническими сбоями и программными уязвимостями. Международный стандарт ISO/IEC 27001 определяет структурированный подход к созданию системы управления информационной безопасностью (СУИБ), основанный на оценке рисков и внедрении адекватных средств защиты \cite{iso27001}.

\ManualSubsection{Основные классы угроз и уязвимостей}

Современные автоматизированные системы подвержены широкому спектру угроз, которые можно классифицировать по источнику и характеру воздействия. Внешние киберугрозы включают в себя атаки на доступность (DDoS-атаки), атаки с использованием вредоносного программного обеспечения (вирусы, ransomware), а также эксплуатацию уязвимостей в сетевых сервисах и прикладном программном обеспечении для несанкционированного доступа к данным. Внутренние угрозы связаны с действиями инсайдеров — как умышленными (саботаж, хищение информации), так и непреднамеренными, возникающими из-за ошибок или неосведомлённости пользователей. Отдельную категорию составляют риски, обусловленные техническими отказами аппаратного обеспечения и системного программного обеспечения, а также уязвимостями, вызванными ошибками проектирования и реализации бизнес-приложений. Растущая сложность IT-ландшафта, включая гибридные и облачные среды, а также распространение устройств интернета вещей (IoT), расширяет периметр атаки и создаёт новые векторы для потенциальных compromise системы \cite{nist}.

\ManualSubsection{Многоуровневая модель защиты и управления доступом}

Эффективная защита автоматизированных систем строится на принципе многоуровневой (эшелонированной) обороны, где меры безопасности реализуются на различных уровнях: физическом, сетевом, операционном, прикладном и данных. На физическом уровне обеспечивается охрана помещений, контроль доступа в дата-центры и защита аппаратных средств. Сетевой уровень защиты включает использование межсетевых экранов (брандмауэров), систем обнаружения и предотвращения вторжений (IDS/IPS), а также сегментацию сети для ограничения распространения возможных атак. На уровне операционных систем и приложений применяется строгое управление доступом на основе принципа наименьших привилегий, регулярное обновление программного обеспечения для устранения известных уязвимостей и использование антивирусных средств. Защита на уровне данных реализуется через криптографические методы — шифрование как данных в состоянии покоя (при хранении), так и данных в движении (при передаче по сети). Центральным элементом модели является система управления идентификацией и доступом (IAM), которая обеспечивает аутентификацию пользователей, авторизацию их действий и учёт (логирование) всех значимых событий для последующего аудита \cite{fstec}.

\ManualSubsection{Обеспечение устойчивости и непрерывности функционирования}

Устойчивость функционирования автоматизированной системы тесно связана с её безопасностью и определяется как способность обеспечивать выполнение критически важных функций в условиях деструктивных воздействий, будь то кибератаки, техногенные аварии или стихийные бедствия. Ключевыми атрибутами устойчивости являются отказоустойчивость, масштабируемость и восстанавливаемость. Отказоустойчивость достигается за счёт архитектурных решений, таких как кластеризация критических сервисов, дублирование сетевых каналов и использование RAID-массивов для хранения данных, что позволяет системе продолжать работу при выходе из строя отдельных компонентов. Важнейшим практическим механизмом обеспечения непрерывности бизнеса является планирование аварийного восстановления (Disaster Recovery Planning, DRP) и разработка планов обеспечения непрерывности бизнеса (Business Continuity Planning, BCP). Эти процессы включают регулярное и автоматизированное резервное копирование данных с хранением копий в географически распределённых локациях, а также проведение регулярных учений по восстановлению для проверки работоспособности процедур \cite{nist}.

\ManualSubsection{Проактивные подходы и тенденции развития}

Современный подход к безопасности смещается от реактивных мер к проактивным и основанным на оценке рисков. Ключевыми элементами такого подхода являются регулярное проведение пентестов (тестов на проникновение) для поиска уязвимостей, статический и динамический анализ безопасности кода (SAST/DAST), а также непрерывный мониторинг безопасности (Security Monitoring) и управление событиями и инцидентами безопасности (SIEM). Актуальной тенденцией является интеграция принципов «безопасности по умолчанию» и «безопасности по проектированию» (Security by Design) в процессы разработки и жизненного цикла программного обеспечения (DevSecOps), что позволяет выявлять и устранять уязвимости на самых ранних стадиях. Развитие технологий искусственного интеллекта и машинного обучения применяется как для создания более совершенных систем защиты, способных обнаруживать сложные аномалии и адаптивные атаки, так и злоумышленниками для автоматизации атак. В условиях ужесточения нормативных требований, таких как ФЗ-152 «О персональных данных» в РФ или GDPR в Европе, соответствие стандартам и проведение регулярных аудитов безопасности становится обязательной практикой для организаций любого масштаба \cite{iso27001}.

\newpage
\refstepcounter{section}
\addcontentsline{toc}{section}{\protect\numberline{\thesection}Примеры применения автоматизированных систем в разных сферах деятельности}
\begin{center}\textbf{\thesection\ ПРИМЕРЫ ПРИМЕНЕНИЯ АВТОМАТИЗИРОВАННЫХ СИСТЕМ В РАЗНЫХ СФЕРАХ ДЕЯТЕЛЬНОСТИ}\end{center}
\vspace*{-1.5em}
\ManualSubsection{Автоматизация в промышленном производстве и логистике}

Промышленное производство является одной из наиболее автоматизированных сфер, где системы управления технологическими процессами (АСУ ТП) интегрируются с корпоративными системами планирования ресурсов предприятия (ERP). Современные автоматизированные системы на производственных предприятиях осуществляют комплексное управление, начиная от контроля состояния оборудования и параметров производственных линий в режиме реального времени и заканчивая оптимизацией логистики поставок сырья и готовой продукции. В отраслях, таких как автомобилестроение и пищевая промышленность, используются роботизированные комплексы для сборки и упаковки, а также системы автоматического контроля качества, основанные на компьютерном зрении. Автоматизация складской логистики с применением систем управления складом (WMS) и роботизированных погрузочно-разгрузочных систем позволяет минимизировать время обработки грузов и снизить количество ошибок \cite{rosstat2024}. Транспортно-логистические компании используют автоматизированные системы для построения маршрутов, спутникового мониторинга подвижного состава и управления цепочками поставок, что существенно повышает их операционную эффективность.

\ManualSubsection{Автоматизированные системы в здравоохранении и медицине}

Сфера здравоохранения активно внедряет автоматизированные информационные системы для решения задач администрирования, клинической практики и управления ресурсами. Медицинские информационные системы (МИС) обеспечивают централизованное хранение и доступ к электронным медицинским картам пациентов, результатам лабораторных и инструментальных исследований, что формирует единую историю болезни и способствует повышению качества диагностики. В области высокотехнологичной медицинской помощи применяются роботизированные хирургические системы, позволяющие проводить операции с минимальной инвазивностью и высокой точностью. Автоматизация процессов в аптечных сетях включает системы управления товарными запасами, контроля сроков годности лекарственных средств и проверки совместимости назначений. Внедрение телемедицинских систем, особенно активизировавшееся в последние годы, предоставляет возможности для дистанционных консультаций, мониторинга хронических заболеваний и организации эффективной работы медицинских учреждений \cite{minzdrav2023}.

\ManualSubsection{Автоматизация в финансовой сфере и государственном управлении}

Финансовый сектор традиционно является лидером по степени автоматизации операционных и аналитических процессов. Банковские автоматизированные системы обеспечивают обработку миллионов транзакций ежедневно через каналы интернет-банкинга, мобильные приложения и платежные терминалы. Ключевыми компонентами являются системы управления взаимоотношениями с клиентами (CRM), платёжные шлюзы, а также сложные системы безопасности и мониторинга мошеннических операций (Fraud Detection), основанные на алгоритмах машинного обучения. Инвестиционные и страховые компании используют автоматизированные аналитические системы и алгоритмические торговые платформы для анализа рыночных данных, оценки рисков и формирования инвестиционных стратегий. В государственном управлении автоматизированные системы лежат в основе цифрового взаимодействия государства с гражданами и бизнесом. Единые порталы государственных услуг предоставляют доступ к широкому спектру административных процедур, а межведомственные системы электронного документооборота и контроля повышают прозрачность и эффективность работы органов власти \cite{hse2023}.

\ManualSubsection{Применение в образовании, транспорте и сельском хозяйстве}

Сфера образования использует автоматизированные системы для управления всеми аспектами учебного процесса. Системы управления обучением (LMS), электронные журналы и деканаты позволяют автоматизировать составление расписаний, ведение успеваемости, организацию самостоятельной работы студентов и доступ к цифровым образовательным ресурсам. В транспортной отрасли автоматизированные системы управления движением обеспечивают безопасность и координацию на авиационном, железнодорожном и городском транспорте. К ним относятся системы управления воздушным движением, автоматизированные системы продажи билетов и контроля пассажиропотока, а также интеллектуальные транспортные системы для управления дорожным движением в городах. В сельском хозяйстве автоматизация проявляется в виде систем точного земледелия, использующих данные геопозиционирования (GPS), датчиков состояния почвы и метеостанций для оптимизации использования ресурсов — воды, семян и удобрений. Автоматизированные комплексы также управляют микроклиматом в теплицах и процессами кормления в животноводстве, что позволяет существенно увеличить производительность и снизить затраты.

\newpage
\section*{}
\vspace*{-2.5em}
\begin{center}\textbf{ЗАКЛЮЧЕНИЕ}\end{center}
\phantomsection
\addcontentsline{toc}{section}{Заключение}

Автоматизированные информационные системы (АИС) представляют собой целостный технологический и организационный комплекс, ставший краеугольным камнем эффективного управления и операционной деятельности во всех сферах современного общества. В ходе проведённого исследования были проанализированы назначение, ключевые характеристики, архитектурные принципы и этапы жизненного цикла АИС, а также определены современные технологические тренды и практические примеры их применения.

Установлено, что фундаментальное назначение АИС заключается в трансформации информационных потоков организации, обеспечивая ускорение процессов, снижение ошибок, систематизацию данных и повышение качества управленческих решений за счёт минимизации человеческого фактора. Качественная архитектура системы, будь то централизованная, децентрализованная или многоуровневая, является основой её надёжности, масштабируемости и безопасности. Управление жизненным циклом АИС, от предпроектного анализа и проектирования до эксплуатации и вывода из эксплуатации, требует системного подхода и адаптации методологий разработки, таких как Agile, для обеспечения соответствия системы динамичным требованиям бизнеса.

Ключевым выводом является конвергенция технологических драйверов, включающих облачные вычисления, аналитику больших данных, искусственный интеллект и интернет вещей, которые переводят АИС на качественно новый уровень, от систем автоматизации рутинных операций к интеллектуальным платформам прогнозной аналитики и адаптивного управления. Внедрение таких систем приносит организациям реальные конкурентные преимущества, выражающиеся в экономии времени и ресурсов, повышении прозрачности процессов и улучшении качества услуг. Однако реализация этого потенциала напрямую зависит от грамотного выбора решений, адекватных задачам организации, и создания комплексной системы информационной безопасности, обеспечивающей конфиденциальность, целостность и доступность данных. Таким образом, автоматизированные информационные системы являются не просто техническим инструментом, а стратегическим активом, определяющим потенциал развития и устойчивости организации в условиях цифровой экономики.

\newpage
\begin{center}\textbf{СПИСОК СОКРАЩЕНИЙ}\end{center}
\phantomsection
\addcontentsline{toc}{section}{Список сокращений}

\vspace{1em}

{
\small
\noindent
\begin{tabular}{|l|p{0.75\textwidth}|}
\hline
\textbf{Сокращение} & \multicolumn{1}{c|}{\textbf{Расшифровка}} \\
\noalign{\hrule height 1.2pt}
\hline
АСУ & Автоматизированная система управления. Класс АИС, предназначенный для управления технологическими, производственными или организационными процессами.\\
\hline
АИС & Автоматизированная информационная система. Организационно-техническая система для автоматизированного выполнения задач управления и обработки данных с участием человека.\\
\hline
ИТ & Информационные технологии. Совокупность методов, процессов и программно-технических средств для работы с информацией.\\
\hline
ИС & Информационная система. Организационно-технический комплекс для сбора, хранения, обработки и распространения информации.\\
\hline
CRM & Управление взаимоотношениями с клиентами (Customer Relationship Management). Система для автоматизации процессов взаимодействия с клиентами.\\
\hline
IoT & Интернет вещей (Internet of Things). Сеть физических объектов, оснащённых встроенными технологиями для взаимодействия друг с другом и с внешней средой.\\
\hline
RPA & Роботизация процессов (Robotic Process Automation). Технология автоматизации рутинных бизнес-процессов с помощью программных роботов.\\
\hline
API & Интерфейс программирования приложений (Application Programming Interface). Набор определений и протоколов для взаимодействия между программными компонентами. \\
\hline
SaaS & Программное обеспечение как услуга (Software as a Service). Модель лицензирования и доставки ПО, при котором оно размещается у провайдера и доступно через интернет.\\
\hline
SDLC & Жизненный цикл разработки программного обеспечения (Software Development Life Cycle). Структурированная последовательность этапов создания программного продукта.\\
\hline
ИИ & Искусственный интеллект. Направление в информатике, занимающееся созданием систем, способных выполнять задачи, требующие человеческого интеллекта.\\
\hline
БД & База данных. Организованная совокупность структурированных данных.\\
\hline
ГИС & Государственная информационная система. Информационная система, создаваемая для осуществления государственных функций и предоставления государственных услуг.\\
\hline
\end{tabular}
}

\clearpage
\nocite{*}
\printbibliography[heading=bibliography]

\end{document}
